\section{Discussion}
The objective was answered affirmatively for conditional equilibrium states: particle-species properties and the cosmological baryon excess permitted separate inference of baryon and antibaryon quantities at a specified literature temperature. The central result at 50 MeV quantifies a nearly symmetric population and a small excess, with 306.1 antibaryons per net unit. As the calculated states cooled, the same ratio reached unity near the 38.2 MeV reference point published by \citet{Yang2022}. The sequence agrees with the general behavior described by \citet{Satarov2013}.

The agreement must be interpreted within the shared conditions. Both approaches connect hadronic equilibrium with cosmological information about the excess per entropy. Fixed inventories and an independent implementation of thermal sums were used here; the kinetic evolution of Satarov et al. was not reproduced, nor was a complete replication of the Yang and Rafelski model performed. The latter impose electric and strangeness neutrality, whereas our calculation maintained $\mu_Q=\mu_S=0$. Their cosmological inputs also differ. Consequently, the difference from 38.2 MeV cannot be attributed exclusively to the catalog or presented as an independent test of early observations.

The simplified chemical condition has a verifiable effect: summing the charges of exported particle species at 50 MeV gave a net charge number of approximately $0.4901\,\Delta n_B$ and a net strangeness number of $-0.03253\,\Delta n_B$. Leptons with zero chemical potential did not compensate this charge. The reconstructed state is therefore a stated thermal approximation; a complete physical cosmological reconstruction would require simultaneous solution of the charge, strangeness, and lepton-sector balances. This limitation remains even though the numerical crossing rounds to the published reference.

The companion evaporation study \citep{GarciaEvap2026} supplies a separate equilibrium control that solves these balances with vanishing net lepton numbers. Under those specific conditions, its crossing differs from this baseline by less than $10^{-8}$ MeV. That local agreement supports the use of the baseline as a numerical reference, while leaving its stated chemical potentials unchanged. The companion kinetic trajectories do not impose all charge balances dynamically.

Comoving quantities clarify the relationship to the present. The excess is conserved under the adopted assumptions, whereas the total equilibrium baryon and antibaryon population decreases. Direct comparison of one early cm$^3$ with one present cm$^3$ would mix population change with dilution by expansion. Entropy normalization permits comparison of the same region and reporting of mean quantities without introducing an undocumented global cosmological volume. H, He, and D abundances serve as external nuclear references from another stage; reconstructing their history requires a subsequent nucleosynthesis network.

A complete historical reconstruction requires rates of inverse production, annihilation, and decay, together with expansion and initial conditions. In a homogeneous framework, an equation of the form
\begin{align}
\frac{dn_{\bar B}}{dt}+3Hn_{\bar B}
&=\langle\sigma v\rangle\bigl(n_B^{\rm eq}n_{\bar B}^{\rm eq}
-n_Bn_{\bar B}\bigr)\label{eq:kinetic}
\end{align}
introduces information absent from an equilibrium table; $H$ is the expansion rate and $\langle\sigma v\rangle$ is the thermally averaged product of cross section and relative velocity. This equation was not integrated in this work. The first event, kinetic endpoint, and cumulative number of regenerated particles were therefore not determined.

The contribution of this work is a reproducible quantitative reference connecting present inputs to early inventories under explicit assumptions, identifying recoverable information and establishing the scope of the published comparison. It does not establish a new observation, an explanation for the origin of the asymmetry, or a new physical mechanism. Numerical convergence and sensitivity to one input do not replace assessment of neutrality, interactions, resonance widths, and the validity of equilibrium.
